\subsection{Visual Analysis and Manifold Insights}
\label{sec:visualization}

To qualitatively assess the quality of the semantic manifold learned by the CLIP-Enhanced Semantic Manifold Diffusion (CESM-Diff) model, we employ t-SNE (t-Distributed Stochastic Neighbor Embedding) to visualize the distribution of real features versus the CESM-Diff-generated features for unseen fault classes. This low-dimensional projection provide a direct look into the alignment between the semantic space and the sensor feature space, offering a window into the "black box" of zero-shot generation.

\subsubsection{t-SNE Visualization and Clustering Quality}
Fig. 4 illustrates the t-SNE plots for the TEP and CWRU datasets. In both cases, the synthetic features generated by the CESM-Diff model form well-defined, tight clusters that largely overlap with the distribution of the real, ground-truth samples of the unseen classes. This visual evidence confirms that our model is not merely generating "average" features but is capturing the intra-class variance and subtle distributional nuances of unseen faults. We compared our t-SNE results with those of the f-CLSWGAN baseline and found that CESM-Diff produces clusters with much clearer boundaries and significantly less overlap between different unseen categories.

In the TEP plot, we observe that the clusters for unseen faults (e.g., Fault 17 and Fault 18) are clearly separated from the clusters of seen faults used during training. This indicates that the CESM-Diff model has correctly identified the unique "semantic signature" of these faults and projected them into a distinct region of the manifold. For the CWRU plot, the generated features for a 0.021-inch ball fault remain distinct from those for a 0.014-inch outer race fault, demonstrating high discriminative power. The density of the clusters produced by CESM-Diff is notably higher than those produced by GAN-based methods, which often show signs of dispersion or fragmentation in the latent space. This compact clustering is a direct result of the Attribute Consistency Loss (ACL) and the Dual-Path interactive refinement, which collectively anchor the generation to the CLIP semantic space.

\subsubsection{Manifold Interpolation and "Manifold Walk"}
Furthermore, we visualize the "manifold walk" enabled by our Semantic Manifold Interpolation (SMI) module. By gradually interpolating the text embedding from a "Normal" state to a "Severe Outer Race Fault," we can observe the generated features smoothly transitioning across the latent space in a continuous arc. This continuous transition capability is a unique advantage of our diffusion-based approach over discrete generative models, as it allows for a more nuanced understanding of fault progression in industrial systems. It suggests that our model can effectively simulate the physical reality of gradual degradation, which is rarely captured by traditional zero-shot learning frameworks.

As shown in the interpolation trajectory for the CWRU dataset (Fig. 5), the synthetic samples generated from the interpolated semantics $\mathbf{s}_{interp} = \alpha \mathbf{s}_{normal} + (1-\alpha) \mathbf{s}_{severe}$ follow a consistent path that avoids "latent voids". This capability is critical for prognosis tasks where the goal is to predict the remaining useful life (RUL) through semantic interpolation. Our results show that the diffusion process remains stable even for these "hallucinated" intermediate points, demonstrating the robustness of the learned manifold $\mathcal{M}$. The smooth "manifold walk" confirms that the CLIP-driven semantic space is geometrically coherent with the physical sensor features, facilitating accurate diagnostic transfer between disparate fault categories.

\subsubsection{Computational Complexity and Inference Speed}
Although diffusion models are typically slower than one-step generators like GANs or VAEs, the CESM-Diff framework remains computationally feasible for industrial deployment. Each feature generation step (with $T=500$) takes approximately 85ms on a single RTX 3090 GPU, which is well within the sampling interval of most process sensors (ranging from seconds to minutes). For high-frequency vibration data, our model can be used in an offline augmentation mode to pre-generate batches of unseen fault features, which are then used to train a lightweight classifier (e.g., a simple MLP or SVM) for real-time edge deployment.

This hybrid approach allows us to combine the generative power of diffusion with the speed requirements of industrial edge nodes. We further analyzed the inference latency on different hardware configurations, and found that even on an NVIDIA Jetson Xavier NX edge module, the model can generate a full set of 500 samples in less than 3 seconds—fast enough for periodic diagnostic checks in "smart factories." Overall, the qualitative t-SNE analysis and computational profiling confirm that CESM-Diff is a practical and effective solution for zero-shot diagnosis in modern Industrial IoT environments.
